\chapter{Conclusion and remarks}
\label{chapter:remarks}

We have discussed in detail the problem of diverse feature selection in linear regression setting. We have established a theoretical result which shows that in a greedy local search (GLS) framework proposed by \citealp{NIPS2012_4689}, using a linear time local search variant based on \citealp{Buchbinder:2012:TLT:2417500.2417835} instead of a rather expensive variant similar to \citealp{Feige:2011:MNS:2079073.2079078}, we can still reach provable optimality guarantee. Unlike the optimality guarantee in \citealp{NIPS2012_4689}, our optimality guarantee does not depend on external parameters to the algorithm.

Also, the present work allows to increase scalability of the entire GLS framework if we use a parallel version of the linear time local search algorithm proposed by \citealp{PanJGBJ14}. This was not possible before due to the sequential nature of the algorithm used in the original paper by \citealp{NIPS2012_4689}. This parallel version selects in each step the exact same element as the one proposed in \citealp{Buchbinder:2012:TLT:2417500.2417835} and therefore the optimality guarantee proposed in the present work holds readily in that case. It would be interesting to see how the usage of the parallel version scales up the size of the problems the framework can handle.

However, from our experiments we have also noticed that the regularizer being non-monotone is a specific case which is hard to observe in many standard data sets. In many cases, although theoretically non-monotone, the regularizers used in our experiments behaved like monotone non-decreasing functions. Therefore, in practice, a standard Forward Regression algorithm suffices to achieve feature subset with sufficiently low error rate. Also, extreme care has to be taken for the regularization constant $\nu$.

Another important thing to note is that in the complexity analysis of the local search algorithms, it is assumed that the function calls are oracle calls in $\mathcal{O}(1)$ which is not the case in practice. Therefore, additional complexity gets added while applying these algorithms in real life specially for larger sized inputs.

To conclude, the usage of submodularity in the feature selection problem makes it interesting from theoretical point of view. It would be interesting to see whether we can design such approach for feature selection in more complex scenario. In particular, in kernel-based feature selection algorithms, such as BAHSIC (see \citealp{SongBBGS07}), which uses Backward Elimination approach to select target features in kernel-mapped space using Hilbert-Schmidt Independence Criterion (HSIC) (see \citealp{Gretton:2005:MSD:2101372.2101382}) as dependency measure, it would be interesting to see if we can achieve some optimality guarantee using submodularity.